

BioSensIA-DC is one computational component of this broader workflow.
It focuses on the structural target-fishing step: given a query
molecule, it ranks candidate protein pockets according to a learned
compatibility score. The prototype is based on DrugCLIP, a contrastive
representation-learning model originally proposed for virtual
screening, and on Uni-Mol-style three-dimensional molecular and pocket
representations \citep{gao23:_drugc,zhou23:_uni_mol}. BioSensIA-DC
uses DrugCLIP as a representation-learning foundation, but extends it
with a new target-fishing workflow, data organization strategy, and
retrieval implementation. In the original
virtual-screening direction, a protein pocket is used as a query to
rank molecules. BioSensIA-DC uses the same molecule--pocket
compatibility idea in the reverse direction: the molecule becomes the
query, and protein pockets become the candidates to be ranked. The
resulting ranked list is intended to support bioreceptor discovery,
not to provide a definitive proof of binding or biosensor performance.


---

\subsection{Virtual screening, target fishing, and representation learning}

